\section{Introdução}

A Ortodontia é, dentre todas as especialidades odontológicas, a que mais rapidamente e organicamente incorporou o paradigma dos gêmeos digitais em sua prática clínica e em sua matriz de pesquisa translacional~\cite{olawade2026advancing}. Essa adoção acelerada não é um fenômeno acidental, mas decorre da natureza inerentemente previsível, mecanicista e iterativa do movimento ortodôntico. Diferentemente de procedimentos cirúrgicos maxilofaciais ou restauradores complexos, cujos resultados dependem de variáveis teciduais estocásticas e de difícil modelagem determinística, a translação e rotação dentárias sob carga ortodôntica seguem princípios biomecânicos e mecanobiológicos rigorosamente estabelecidos. Tais princípios podem ser capturados, parametrizados e simulados computacionalmente com um nível de fidelidade que tangencia a precisão \textit{in vivo}~\cite{zhang2024recursive}.

A capacidade de orquestrar a movimentação dentária em um ambiente virtual holístico antes de aplicar vetores de força clinicamente revolucionou não apenas o planejamento, mas a própria ontologia do tratamento ortodôntico. Historicamente dependente de heurísticas e de ajustes reativos (o chamado \textit{wire-bending} empírico), o clínico contemporâneo opera com capacidades preditivas, podendo visualizar, quantificar, ajustar e otimizar cada incremento espacial de deslocamento em um modelo virtual isomorfo ao paciente -- seu gêmeo digital. Essa transição paradigmática reduz de forma contundente o tempo de cadeira, a incidência de reabsorções radiculares iatrogênicas e a imprevisibilidade outrora associada às abordagens puramente analógicas~\cite{duggal2026exploring}.

Neste cenário de confluência entre a biomecânica clássica e a ciência da computação avançada~\cite{mdpi2025convergence}, este capítulo dedica-se a explorar exaustivamente os alicerces teóricos que tornam a simulação ortodôntica um domínio de excelência para os gêmeos digitais. \destaque{Examinaremos de forma detida a aplicação do Método dos Elementos Finitos (FEM) na predição de centros de resistência e eixos de rotação}, dissecaremos o estado da arte das réplicas virtuais aplicadas aos sistemas de alinhadores transparentes (\textit{Clear Aligner Therapy} -- CAT) e elucidaremos os gargalos computacionais e biológicos que ainda impõem limites à adoção irrestrita destas tecnologias~\cite{li2025impacts, li2024aligner}. Como ficará patente nas seções subsequentes, a Ortodontia calcada em alinhadores e ancoragem esquelética guiada constitui, inequivocamente, o ecossistema de caso de uso mais maduro, metodologicamente escrutinado e clinicamente validado de gêmeos digitais na Odontologia de precisão contemporânea~\cite{techsoc2025precision}.
